\chapter{INTRODUÇÃO}

Métodos automáticos permitem acompanhar concentrações de material particulado em escalas temporais menores que as usualmente obtidas pela coleta e pesagem de filtros. Essa resolução favorece o acompanhamento de episódios e a investigação da variabilidade das concentrações. Entretanto, a disponibilidade de dados contínuos não garante sua intercambialidade com a gravimetria. Os métodos respondem a propriedades e condições de amostragem diferentes, e parte da massa coletada pode ser modificada antes da determinação final \parencite{rsl204,rsl201,rsl045}.

Entre os métodos automáticos convencionais estão a microbalança de elemento oscilante cônico, do inglês \textit{Tapered Element Oscillating Microbalance} (TEOM), suas modificações, os monitores por atenuação beta e os instrumentos baseados em espalhamento de luz. Aquecimento, controle de umidade, vazão, entrada seletiva e algoritmos alteram o resultado. Na inferência da massa a partir de propriedades ópticas ou da distribuição de tamanho das partículas, também são relevantes composição, densidade e forma do aerossol \parencite{rsl288,rsl283,rsl376}.

O comparador gravimétrico tampouco é imune a artefatos. Coleta, armazenamento e condicionamento podem favorecer perdas semivoláteis, retenção de água e diferenças entre filtros ou amostradores. A discrepância pode, portanto, resultar de ambos os lados da comparação. Interpretá-la exige identificar a massa operacionalmente determinada e os procedimentos utilizados \parencite{rsl202,rsl101,rsl105}.

Essa questão é relevante para redes que combinam equipamentos e períodos de monitoramento diferentes. Mudanças instrumentais ou de correção podem alterar a distribuição das concentrações e afetar análises de tendência e representatividade temporal. É necessário compreender em quais condições os métodos se aproximam e quais fatores restringem a transferência de resultados entre locais.

O objetivo desta revisão foi avaliar a comparabilidade entre métodos automáticos convencionais e gravimétricos para \MPdez{} e \MPdoiscinco{} em comparações externas, colocalizadas e concomitantes, sintetizando resultados de desvio e identificando fatores ambientais, operacionais e instrumentais associados às diferenças. Sensores de baixo custo foram excluídos como métodos de interesse. Nos estudos norte-americanos, distinguem-se o método federal de referência, \textit{Federal Reference Method} (FRM), e os métodos federais equivalentes, \textit{Federal Equivalent Methods} (FEM); essas designações não dispensam o exame do desenho de cada comparação.
